\documentclass[12pt,a4paper]{article}

% --- Essential Packages ---
\usepackage[utf8]{inputenc}
\usepackage[T1]{fontenc}
\usepackage{amsmath, amssymb, amstext} % Math support [cite: 98]
\usepackage{geometry} % For margins
\usepackage{booktabs} % For professional tables [cite: 34, 43, 98]
\usepackage{caption} % For table/figure captions
\usepackage{hyperref} % For clickable links and references
\usepackage{setspace} % For line spacing

% --- Layout Settings ---
\geometry{margin=1in}
\onehalfspacing

\title{Technical Appendix: The Fama-French (2010) Bootstrap Methodology}
\author{Ömer}
\date{March 2026}

\begin{document}

\maketitle

\section{Methodological Framework}
The central objective of the bootstrap analysis is to separate managerial skill from statistical luck within the cross-section of mutual fund alphas[cite: 97, 98]. In a large universe of funds, extreme performance—both positive and negative—can arise solely through the distribution of luck. This methodology generates a ``zero-skill'' benchmark to serve as a yardstick for evaluating realized performance[cite: 98].

\section{The Step-by-Step Bootstrap Process}

\subsection{Establishing the Null Hypothesis}
To test for luck, we first create a synthetic universe where no manager possesses forecasting talent. We calculate the realized alpha ($\hat{\alpha}_i$) for each fund $i$ using the Carhart four-factor model. We then subtract this realized alpha from the fund’s monthly return series:
\begin{equation}
    R_{i,t}^{Adj} = R_{i,t} - \hat{\alpha}_i
\end{equation}
By construction, every fund in this bootstrap universe now has a true population alpha of exactly zero[cite: 98].

\subsection{Resampling and Cross-Sectional Correlation}
We draw $T$ months with replacement from the sample period (1995--2023). Crucially, for any given bootstrap run, the \textit{same} sequence of random months is applied to every fund. This preserves the cross-sectional correlation of returns—ensuring that common market shocks affecting the industry are maintained in the simulated data.

\subsection{Generating the Simulated Distribution}
We execute 10,000 bootstrap runs. In each run, we perform the four-factor regression for every fund using the resampled data and record the $t$-statistic of the alpha ($t(\alpha)$). Within each run, funds are ranked from the 1st percentile (worst) to the 99th percentile (best) based on their $t(\alpha)$[cite: 98].

\subsection{Probability of Luck}
The probability of luck is determined by comparing the realized $t(\alpha)$ at a specific percentile to the average simulated $t(\alpha)$ at that same percentile. For a result to be considered statistically significant skill (or significant incompetence), the actual $t$-statistic must fall outside the 95\% confidence interval of the luck-only distribution[cite: 98].

\section{Advanced Metrics: $\pi_0$ Estimation}
Following \textcite{FamaFrench2010}, we estimate the proportion of the entire fund population that has a true alpha of exactly zero, denoted as $\pi_0$[cite: 82, 100]. This is calculated using the density of $p$-values in the middle of the distribution ($\lambda = 0.5$):
\begin{equation}
    \hat{\pi}_0 = \frac{\#\{p_i > \lambda\}}{N \times (1 - \lambda)}
\end{equation}
For this sample, we estimate $\hat{\pi}_0 = 76.8\%$, suggesting that roughly 77\% of the mutual fund industry possesses zero true talent[cite: 82, 99].

\section{Summary of Findings}
\begin{itemize}
    \item \textbf{The Left Tail:} Realized performance at the 1st percentile ($t = -3.518$) is significantly worse than the simulated luck mean of $-2.368$, with only a 0.3\% probability of being luck-driven[cite: 98]. This confirms the existence of persistently incompetent managers or excessive fee structures.
    \item \textbf{The Right Tail:} Performance at the 99th percentile ($t = 2.431$) is statistically indistinguishable from the simulated luck mean of $2.340$, as 63.4\% of zero-skill runs produced lower values[cite: 98].
\end{itemize}

\end{document}